\documentclass[10pt,a4paper]{article}
\usepackage[a4paper,margin=14mm,top=12mm,bottom=14mm]{geometry}
\usepackage{fontspec}
\setmainfont{Liberation Sans}
\usepackage{polyglossia}
\setdefaultlanguage{french}
\usepackage{amsmath}
\usepackage{xcolor}
\usepackage{graphicx}
\usepackage[most]{tcolorbox}
\usepackage{tabularx,booktabs,array,multirow}
\usepackage{enumitem}
\usepackage{fancyhdr}
\usepackage{lastpage}
\usepackage[hidelinks]{hyperref}
\definecolor{uimm}{HTML}{005677}
\definecolor{cyan}{HTML}{13A2DD}
\definecolor{fond}{HTML}{E6F2FA}
\definecolor{ok}{HTML}{2E7D32}
\definecolor{ko}{HTML}{B3261E}
\definecolor{violet}{HTML}{7B1FA2}
\setlength{\parindent}{0pt}
\setlength{\parskip}{3pt}
\setlist{nosep}
\setlist[itemize]{leftmargin=13pt,label=\textcolor{cyan}{\textbullet}}
\newcommand{\titre}[1]{\vspace{5pt}{\color{uimm}\large\bfseries #1}\par\vspace{-3pt}{\color{cyan}\rule{\linewidth}{1.2pt}}\par\vspace{0pt}}
\tcbset{boite/.style={colback=fond,colframe=uimm,boxrule=0.6pt,arc=2pt,left=6pt,right=6pt,top=4pt,bottom=4pt},
        defi/.style={colback=violet!6,colframe=violet,boxrule=0.8pt,arc=2pt,left=6pt,right=6pt,top=4pt,bottom=4pt,
                     title={\bfseries Défi},coltitle=white,colbacktitle=violet,fonttitle=\small}}
\newcommand{\bandeau}[3]{%
\begin{tcolorbox}[colback=uimm,colframe=uimm,arc=2pt,left=4pt,right=8pt,top=4pt,bottom=4pt]
\begin{minipage}[c]{18mm}\setlength{\fboxsep}{1.2pt}\fcolorbox{white}{white}{\includegraphics[height=16.5mm]{logo_uimm.jpg}}\end{minipage}\hspace{4mm}%
\begin{minipage}[c]{\dimexpr\linewidth-22mm\relax}\color{white}
{\small #1}\\[1pt]{\LARGE\bfseries #2}\\[2pt]{\large #3}\hfill{\small CQPM Technicien qualité}
\end{minipage}
\end{tcolorbox}}
\newcommand{\reponse}[1][1]{\par\vspace{2pt}{\color{gray!60}\foreach \i in {1,...,#1}{\rule{\linewidth}{0.3pt}\par\vspace{9pt}}}}
\usepackage{pgffor}
\pagestyle{fancy}\fancyhf{}
\renewcommand{\headrulewidth}{0pt}
\newcommand{\pied}[1]{\fancyfoot[L]{\footnotesize\color{uimm}#1}\fancyfoot[R]{\footnotesize\color{uimm}Page \thepage{} sur \pageref{LastPage}}}

\usepackage{tikz}
\usepackage{longtable}
\renewcommand{\bandeau}[3]{%
\begin{tcolorbox}[colback=uimm,colframe=uimm,arc=2pt,left=4pt,right=8pt,top=4pt,bottom=4pt]
\begin{minipage}[c]{18mm}\setlength{\fboxsep}{1.2pt}\fcolorbox{white}{white}{\includegraphics[height=16.5mm]{logo_uimm.jpg}}\end{minipage}\hspace{4mm}%
\begin{minipage}[c]{\dimexpr\linewidth-22mm\relax}\color{white}
{\small #1}\\[1pt]{\LARGE\bfseries #2}\\[2pt]{\large #3}
\end{minipage}
\end{tcolorbox}}
% en-tête d'une fiche : code, titre, question, bloc, durée, organisation, niveau
\newcommand{\entete}[7]{%
\begin{tcolorbox}[colback=uimm,colframe=uimm,arc=2pt,left=4pt,right=6pt,top=4pt,bottom=4pt]
\begin{minipage}[c]{18mm}\setlength{\fboxsep}{1.2pt}\fcolorbox{white}{white}{\includegraphics[height=16.5mm]{logo_uimm.jpg}}\end{minipage}\hspace{3mm}%
\begin{minipage}[c]{\dimexpr\linewidth-58mm\relax}\color{white}\raggedright
{\small Base SPC · #4}\\[1pt]{\Large\bfseries #1 · #2}\\[2pt]{\itshape #3}
\end{minipage}\hfill
\begin{minipage}[c]{33mm}\color{white}\small\raggedleft
Durée : \textbf{#5}\\ \textbf{#6}\\ Niveau #7 sur 3
\end{minipage}
\end{tcolorbox}}
\newcommand{\tagE}{{\setlength{\fboxsep}{1.5pt}\colorbox{cyan}{\color{white}\scriptsize\bfseries Ellistat}}}
\newcommand{\tagW}{{\setlength{\fboxsep}{1.5pt}\colorbox{uimm}{\color{white}\scriptsize\bfseries Page web}}}
\tcbset{pred/.style={colback=orange!6,colframe=orange!80!black,boxrule=0.6pt,arc=2pt,left=6pt,right=6pt,top=3pt,bottom=1pt,
                    title={\bfseries Je prédis, avant de commencer},coltitle=white,colbacktitle=orange!80!black,fonttitle=\small},
        retenir/.style={colback=fond,colframe=uimm,boxrule=0.6pt,arc=2pt,left=6pt,right=6pt,top=3pt,bottom=3pt,
                    title={\bfseries À retenir},coltitle=white,colbacktitle=uimm,fonttitle=\small},
        plus/.style={colback=violet!6,colframe=violet,boxrule=0.6pt,arc=2pt,left=6pt,right=6pt,top=3pt,bottom=3pt,
                    title={\bfseries Pour aller plus loin},coltitle=white,colbacktitle=violet,fonttitle=\small},
        act/.style={colback=uimm!10,colframe=uimm,boxrule=0.5pt,arc=2pt,left=5pt,right=5pt,top=2pt,bottom=2pt,before skip=7pt,after skip=3pt},
        verif/.style={colback=orange!6,colframe=orange!80!black,boxrule=0.6pt,arc=2pt,left=6pt,right=6pt,top=3pt,bottom=3pt,
                    title={\bfseries À vérifier sur poste avant la première séance},coltitle=white,colbacktitle=orange!80!black,fonttitle=\small}}
\newcommand{\fichier}[1]{\texttt{\small #1}}
\newcommand{\menu}[1]{\textsf{\textbf{#1}}}
\newcommand{\champ}[1]{\rule{#1}{0.3pt}}

\usepackage{colortbl}
\pied{Base SPC · Fiche formateur · Gage R\&R}
\newcommand{\stitre}[1]{\par\vspace{1pt}{\color{uimm}\bfseries\large #1}\par\vspace{-7pt}{\color{cyan}\rule{\linewidth}{0.9pt}}\par\vspace{-3pt}}
\newcommand{\CM}[1]{\mathit{CM}_{#1}}
\newcommand{\SC}[1]{\mathit{SC}_{#1}}
\begin{document}
\thispagestyle{fancy}
\fontsize{9.2}{11}\selectfont
\bandeau{Base SPC · Fiche formateur}{Gage R\&R : rappel des formules}{Du modèle au \%GRR, avec un exemple à reproduire dans Ellistat}
\vspace{2pt}

\stitre{1. Le modèle}
$n$ pièces, $k$ opérateurs, $r$ répétitions (exemple Bosch : $n=10$, $k=3$, $r=2$). La mesure de la pièce $i$ par l'opérateur $j$, à la répétition $l$ :
\begin{center}\vspace{-3pt}$x_{ijl}=\mu+P_i+O_j+(PO)_{ij}+\varepsilon_{ijl}$\vspace{-2pt}\end{center}
Les quatre termes sont aléatoires, de variances $\sigma_P^2$ (pièces), $\sigma_O^2$ (opérateurs), $\sigma_{PO}^2$ (interaction pièce $\times$ opérateur) et $\sigma^2$ (répétition). Ces variances s'additionnent, pas les écarts types.\par\vspace{3pt}
\renewcommand{\arraystretch}{1.08}
\begin{tabularx}{\linewidth}{|>{\bfseries}p{37mm}|>{\raggedright\arraybackslash}X|}\hline
\rowcolor{uimm}\color{white}\textbf{Nom} & \color{white}\textbf{Définition}\\\hline
EV, répétabilité & $\sigma$, l'écart type des mesures répétées (même pièce, même opérateur)\\\hline
AV, reproductibilité & $\sigma_O$, l'écart type entre opérateurs (avec $\sigma_{PO}$ si l'interaction n'est pas regroupée)\\\hline
PV, pièces & $\sigma_P$, l'écart type entre pièces\\\hline
GRR & $\sqrt{\mathrm{EV}^2+\mathrm{AV}^2}$ \quad (le «\,Gage R\&R\,», en unité de mesure)\\\hline
TV, total & $\sqrt{\mathrm{GRR}^2+\mathrm{PV}^2}$\\\hline
\end{tabularx}

\stitre{2. Le tableau d'analyse de variance}
Moyennes : $\bar x_{i\cdot\cdot}$ de la pièce $i$ (sur $kr$ mesures), $\bar x_{\cdot j\cdot}$ de l'opérateur $j$ (sur $nr$), $\bar x_{ij\cdot}$ de la cellule (sur $r$), $\bar x_{\cdot\cdot\cdot}$ générale.\par\vspace{3pt}
\renewcommand{\arraystretch}{1.28}
\begin{tabularx}{\linewidth}{|>{\bfseries\raggedright\arraybackslash}p{24mm}|c|>{\raggedright\arraybackslash}X|c|}\hline
\rowcolor{uimm}\color{white}\textbf{Source} & \color{white}\textbf{ddl} & \color{white}\textbf{Somme des carrés $\SC{}$} & \color{white}\textbf{Espérance de $\CM{}=\SC{}/\mathrm{ddl}$}\\\hline
Pièces & $n-1$ & $kr\sum_i(\bar x_{i\cdot\cdot}-\bar x_{\cdot\cdot\cdot})^2$ & $\sigma^2+r\sigma_{PO}^2+kr\sigma_P^2$\\\hline
Opérateurs & $k-1$ & $nr\sum_j(\bar x_{\cdot j\cdot}-\bar x_{\cdot\cdot\cdot})^2$ & $\sigma^2+r\sigma_{PO}^2+nr\sigma_O^2$\\\hline
Interaction & $(n-1)(k-1)$ & $r\sum_{i,j}(\bar x_{ij\cdot}-\bar x_{i\cdot\cdot}-\bar x_{\cdot j\cdot}+\bar x_{\cdot\cdot\cdot})^2$ & $\sigma^2+r\sigma_{PO}^2$\\\hline
Répétabilité & $nk(r-1)$ & $\sum_{i,j,l}(x_{ijl}-\bar x_{ij\cdot})^2$ & $\sigma^2$\\\hline
Total & $nkr-1$ & $\sum_{i,j,l}(x_{ijl}-\bar x_{\cdot\cdot\cdot})^2$ & \\\hline
\end{tabularx}
\par\smallskip
La dernière colonne donne les estimateurs : on résout ces équations pour $\sigma^2$, $\sigma_O^2$, $\sigma_P^2$.

\stitre{3. Du tableau au \%GRR}
\begin{enumerate}[leftmargin=14pt,itemsep=2pt,topsep=2pt]
\item \textbf{Interaction.} $F=\CM{PO}/\CM{E}$, avec $(n-1)(k-1)$ et $nk(r-1)$ ddl. Si $p>0{,}05$ (non significative), on la regroupe avec la répétabilité : $\CM{E}'=\dfrac{\SC{PO}+\SC{E}}{\mathrm{ddl}_{PO}+\mathrm{ddl}_{E}}$. C'est ce que fait Ellistat (ligne «\,interaction 0\,»).
\item \textbf{Écarts types.} $\mathrm{EV}=\sqrt{\CM{E}'}$ \quad $\mathrm{AV}=\sqrt{\dfrac{\CM{O}-\CM{E}'}{nr}}$ \quad $\mathrm{PV}=\sqrt{\dfrac{\CM{P}-\CM{E}'}{kr}}$. Sous la racine, un nombre négatif est ramené à 0.
\item \textbf{Combiner.} $\mathrm{GRR}=\sqrt{\mathrm{EV}^2+\mathrm{AV}^2}$ \quad $\mathrm{TV}=\sqrt{\mathrm{GRR}^2+\mathrm{PV}^2}$.
\item \textbf{Comparer à la tolérance $T$.} $\%\mathrm{GRR}=\dfrac{6\,\mathrm{GRR}}{T}$ (idem \%EV, \%AV) \quad $C_{pc}=\dfrac{T}{6\,\mathrm{GRR}}=\dfrac{100}{\%\mathrm{GRR}}$ \quad $ndc=\sqrt2\,\dfrac{\mathrm{PV}}{\mathrm{GRR}}$ \quad $\dfrac{\mathrm{GRR}}{\mathrm{TV}}$ (part de la variation totale).
\item \textbf{Sans regroupement} (interaction significative) : $\mathrm{EV}=\sqrt{\CM{E}}$, $\mathrm{AV}=\sqrt{(\CM{O}-\CM{PO})/(nr)}$, $\mathrm{INT}=\sqrt{(\CM{PO}-\CM{E})/r}$, $\mathrm{GRR}=\sqrt{\mathrm{EV}^2+\mathrm{AV}^2+\mathrm{INT}^2}$.
\end{enumerate}

\stitre{4. Lire le résultat}
\renewcommand{\arraystretch}{1.1}
\begin{tabularx}{\linewidth}{|>{\bfseries\raggedright\arraybackslash}p{37mm}|>{\raggedright\arraybackslash}X|}\hline
Le facteur 6 & $2K$ avec $K=3$ (AIAG, réglage d'Ellistat) : la mesure est comptée sur $\pm 3$ écarts types. Avec $K=2$ (VDA) le facteur devient 4.\\\hline
Livret Bosch, AIAG & $\%\mathrm{GRR}<10$ \% capable ; de 10 à 30 \% capable sous conditions ; plus de 30 \% non capable. Et $ndc\ge 5$.\\\hline
Ellistat & Jugement par le $C_{pc}$ : rouge sous 3 ($\%\mathrm{GRR}>33$ \%), orange de 3 à 4, vert au-dessus de 4 ($<25$ \%). Les deux jugements divergent entre 10 et 33 \% : à faire discuter.\\\hline
Tolérance minimale & Pour 10 \% : $T_{\min}=60\,\mathrm{GRR}$. Ellistat affiche $24\,\mathrm{GRR}$ ($C_{pc}=4$).\\\hline
\end{tabularx}

\newpage
\bandeau{Base SPC · Fiche formateur}{Exemple à reproduire dans Ellistat}{Livret Bosch p.\,14 : 10 pièces, 3 opérateurs, 2 répétitions}
\vspace{3pt}
\begin{minipage}[t]{0.535\linewidth}\vspace{0pt}
\stitre{Les données (mm)}
\renewcommand{\arraystretch}{1.0}\setlength{\tabcolsep}{3.5pt}
{\small\begin{tabular}{|c|c|c|c|c|c|c|>{\columncolor{fond}}c|}\hline
\rowcolor{uimm}\color{white}\textbf{Pièce} & \color{white}A1 & \color{white}A2 & \color{white}B1 & \color{white}B2 & \color{white}C1 & \color{white}C2 & \color{white}$\bar x_{i\cdot\cdot}$\\\hline
1 & 6,029 & 6,030 & 6,033 & 6,032 & 6,031 & 6,030 & 6,0308\\
2 & 6,019 & 6,020 & 6,020 & 6,019 & 6,020 & 6,020 & 6,0197\\
3 & 6,004 & 6,003 & 6,007 & 6,007 & 6,010 & 6,006 & 6,0062\\
4 & 5,982 & 5,982 & 5,985 & 5,986 & 5,984 & 5,984 & 5,9838\\
5 & 6,009 & 6,009 & 6,014 & 6,014 & 6,015 & 6,014 & 6,0125\\
6 & 5,971 & 5,972 & 5,973 & 5,972 & 5,975 & 5,974 & 5,9728\\
7 & 5,995 & 5,997 & 5,997 & 5,996 & 5,995 & 5,994 & 5,9957\\
8 & 6,014 & 6,018 & 6,019 & 6,015 & 6,016 & 6,015 & 6,0162\\
9 & 5,985 & 5,987 & 5,987 & 5,986 & 5,987 & 5,986 & 5,9863\\
10 & 6,024 & 6,028 & 6,029 & 6,025 & 6,026 & 6,025 & 6,0262\\\hline
\end{tabular}}
\par\smallskip
{\small Moyennes des opérateurs $\bar x_{\cdot j\cdot}$ : A 6,00390 ; B 6,00580 ; C 6,00535. Moyenne générale 6,005017.\par
Tolérance $T=0{,}060$ (de 5,97 à 6,03).\par
Fichier : \fichier{SPC\textbackslash Donnees\_Ellistat\textbackslash RR\_Bosch\_p14.xlsx}}
\end{minipage}\hfill
\begin{minipage}[t]{0.45\linewidth}\vspace{0pt}
\stitre{Le tableau (en $\mu$m$^2$)}
\renewcommand{\arraystretch}{1.12}\setlength{\tabcolsep}{4pt}
{\small\begin{tabular}{|l|c|r|r|}\hline
\rowcolor{uimm}\color{white}Source & \color{white}ddl & \color{white}$\SC{}$ & \color{white}$\CM{}$\\\hline
Pièces & 9 & 20\,586,48 & 2\,287,39\\\hline
Opérateurs & 2 & 39,43 & 19,717\\\hline
Interaction & 18 & 60,57 & 3,365\\\hline
Répétabilité & 30 & 52,50 & 1,750\\\hline
Total & 59 & 20\,738,98 & \\\hline
\end{tabular}}
\par\smallskip
{\small 1 $\mu$m $=0{,}001$ mm : on calcule en $\mu$m pour éviter les puissances de 10.}
\stitre{Les étapes}
{\small\begin{enumerate}[leftmargin=12pt,itemsep=1.5pt,topsep=1pt]
\item $F=3{,}365/1{,}750=1{,}92$. Seuil $F_{0{,}95}(18\,;30)=1{,}96$, $p=0{,}055$ : non significative, donc regroupée. $\CM{E}'=\dfrac{60{,}57+52{,}50}{48}=2{,}3556$.
\item $\mathrm{EV}=\sqrt{2{,}3556}=1{,}5348$. $\mathrm{AV}=\sqrt{\dfrac{19{,}717-2{,}3556}{20}}=0{,}9317$. $\mathrm{PV}=\sqrt{\dfrac{2287{,}39-2{,}3556}{6}}=19{,}515$.
\item $\mathrm{GRR}=\sqrt{2{,}3556+0{,}8681}=1{,}7954$. $\mathrm{TV}=\sqrt{3{,}2237+380{,}84}=19{,}598$.
\item $\%\mathrm{GRR}=6\times1{,}7954/60=17{,}95$ \%.
\end{enumerate}}
\end{minipage}

\stitre{Ce que tu dois lire dans Ellistat (K = 3), écran vu le 14/11/2026}
\renewcommand{\arraystretch}{1.08}\setlength{\tabcolsep}{4pt}
{\small\begin{tabularx}{\linewidth}{|>{\raggedright\arraybackslash}X|c|c|>{\raggedright\arraybackslash}X|c|c|}\hline
\rowcolor{uimm}\color{white}Ellistat & \color{white}Mon calcul & \color{white}Écran & \color{white}Ellistat & \color{white}Mon calcul & \color{white}Écran\\\hline
Répétabilité (EV) & 0,001535 & 0,001535 & Écart type global (TV) & 0,01960 & 0,01960\\\hline
Reproductibilité (AV) & 0,0009317 & 0,0009317 & Nombre de catégories & 15,37 & 15,4\\\hline
Instrument (GRR) & 0,001795 & 0,001795 & GRR \% et $C_{pc}$ & 17,95 \% ; 5,57 & 18,0 \% ; 5,570\\\hline
Pièces (PV) & 0,01952 & 0,01952 & GRR/TV («\,production\,») & 9,16 \% & 9,16 \%\\\hline
Interaction & regroupée & 0 & Tolérance minimale & 0,04309 & 0,04309\\\hline
\end{tabularx}}

\stitre{Dans Ellistat}
\begin{enumerate}[leftmargin=14pt,itemsep=1pt,topsep=1pt]
\item \menu{Gage R\&R}, \menu{Mesures}, \menu{Ajouter un gage R\&R}. Nom ; Produits 10 ; Opérateurs 3 ; \textbf{Répétitions 2} (3 par défaut) ; méthode ANAVAR ; $K=3$ (AIAG).
\item Piège : les tolérances saisies à la création ne sont pas conservées. Ouvrir \menu{Voir Plus} à côté de \menu{Informations du gage} et saisir 5,97 et 6,03.
\item Onglet \menu{Mesure} : un tableau par opérateur, deux colonnes. Coller A1 A2, puis B1 B2, puis C1 C2, sans en-tête. Onglet \menu{Analyse} (bouton $>$) pour comparer.
\end{enumerate}

\stitre{Application : G2 à 4 opérateurs (arbres Ø\,48, jeu simulé)}
{\small Fichier \fichier{RR\_arbres\_diametre48\_4operateurs\_simule.xlsx}, 10 pièces, 4 opérateurs, 2 répétitions. Mon calcul : EV 0,00542 ; AV 0,00315 ; GRR 0,00627 ; PV 0,1057 ; $ndc=23{,}9$. Interaction $F=0{,}76$, $p=0{,}77$ : regroupée. Opérateur $F=7{,}8$ pour un seuil de 2,74 : significatif, mais AV pèse 25 \% de GRR$^2$.\par\smallskip
\renewcommand{\arraystretch}{1.05}
\begin{tabularx}{\linewidth}{|c|c|c|>{\raggedright\arraybackslash}X|}\hline
\rowcolor{uimm}\color{white}$T$ (mm) & \color{white}\%GRR & \color{white}$C_{pc}$ & \color{white}Verdict du livret\\\hline
0,4 & 9,40 \% & 10,6 & \cellcolor{ok!25}capable\\\hline
0,2 & 18,8 \% & 5,3 & \cellcolor{orange!25}capable sous conditions\\\hline
0,1 & 37,6 \% & 2,7 & \cellcolor{ko!25}non capable\\\hline
0,05 & 75,2 \% & 1,3 & \cellcolor{ko!25}non capable\\\hline
\end{tabularx}}

\begin{tcolorbox}[verif,title={\bfseries Pièges et limites},top=2pt,bottom=2pt]
{\small \textbf{Interaction à la limite.} Bosch : $p=0{,}055$. Sans regroupement, GRR $=1{,}8371$ $\mu$m soit 18,4 \% (calcul, non vu à l'écran). Le verdict ne change pas, mais le regroupement est un choix. \textbf{Significatif n'est pas important.} Opérateur Bosch : $F=8{,}37$ pour un seuil de 3,19, mais AV vaut 9,3 \% de $T$. \textbf{Non vu à l'écran :} Ellistat à 4 opérateurs, donc les valeurs G2 sont des calculs. Elles sont à comparer, et un écart viendrait du réglage ($K$, regroupement) plutôt que de la méthode. \textbf{Sources :} livret Bosch, \emph{Technical Statistics}, Booklet 10 ; AIAG, \emph{MSA}, 4\textsuperscript{e} éd.}
\end{tcolorbox}
\end{document}
